%%%%%%%%%%%%%%%%%%%%%%%%%%%%%%%%%%
% Página de resumen del proyecto %
%%%%%%%%%%%%%%%%%%%%%%%%%%%%%%%%%%
% !TEX root = A0.MiTFG.tex

\pagestyle{fancy}
\renewcommand{\headrulewidth}{0pt}
\addstarredchapter{Resumen}

\begin{center}
	\scshape
	E.T.S. de Ingeniería de Telecomunicación, Universidad de Málaga
\end{center}

\bigskip

\begin{center}
	\Large \scshape
	\textbf{\tfgtitlename}
\end{center}

\bigskip \bigskip \bigskip

\begin{minipage}{\textwidth}

\textbf{Autor:} \tfgauthorname

\medskip

\textbf{Tutor:} \tfgtutorname

\medskip

\textbf{Departamento:} Tecnología Electrónica 

\medskip

\textbf{Titulación:} Grado en Ingeniería de sistemas electrónicos 

\medskip

\textbf{Palabras clave:} RISC-V, FPGA, arquitectura de CPU, hardware de código abierto, diseño digital, simulación, HDL, Verilog, procesador embebido, sistema en chip, caché, pipeline, validación funcional, desarrollo de hardware.

\bigskip \bigskip


\end{minipage}

\begin{center}
	\textbf{Resumen}
\end{center}

El presente proyecto tiene como objetivo el diseño de un procesador basado en la arquitectura RISC-V, implementando específicamente el conjunto de instrucciones RV32I. Para garantizar su correcto funcionamiento, se emplea un método de verificación fundamentado en la generación y comparación de firmas. Finalmente, el procesador se implementa en una plataforma FPGA con el fin de validar experimentalmente su desempeño y comprobar la fidelidad respecto a la especificación arquitectónica.


\blankpage
